\documentclass[14pt]{extarticle}
\usepackage{natbib}

\usepackage{xcolor}
\usepackage{amsmath}
\newcommand{\tuple}[1]{ \langle #1 \rangle }
%\usepackage{automata}
\usepackage{times}
\usepackage{ltablex}
\usepackage{tasks}
\usepackage{threeparttable}
\usepackage{booktabs}
\usepackage{xcolor} % For custom colors
\usepackage{tikz} % For styling enumerate numbers
\usepackage{tcolorbox} % For colored box styling
\usepackage{amsmath, amsfonts, amssymb, amsthm} % Math related
\usepackage{natbib}
\usepackage{fontspec}
\usepackage{luatexja}
\usepackage[mathscr]{euscript}

%%%%%% Template
\usepackage{hyperref}
\hypersetup{colorlinks=true,allcolors=blue}

\usepackage{vmargin}
\setpapersize{USletter}
\setmarginsrb{1.0in}{1.0in}{1.0in}{0.6in}{0pt}{0pt}{0pt}{0.4in}

% HOW TO USE THE ABOVE:
%\setmarginsrb{leftmargin}{topmargin}{rightmargin}{bottommargin}{headheight}{headsep}{footheight}{footskip}
%\raggedbottom
% paragraphs indent & skip:
\parindent  0.3cm
\parskip    -0.01cm

\usepackage{tikz}
\usetikzlibrary{backgrounds}

% hyphenation:
% \hyphenpenalty=10000 % no hyphen
% \exhyphenpenalty=10000 % no hyphen
\sloppy

% notes-style paragraph spacing and indentation:
\usepackage{parskip}
\setlength{\parindent}{0cm}

% let derivations break across pages
\allowdisplaybreaks

\newcommand{\orange}[1]{\textcolor{orange}{#1}}
\newcommand{\blue}[1]{\textcolor{blue}{#1}}
\newcommand{\red}[1]{\textcolor{red}{#1}}
\newcommand{\freq}[1]{{\bf \sf F}(#1)}
\newcommand{\datafreq}[2]{{{\bf \sf F}_{#1}(#2)}}

\def\qqquad{\quad\qquad}
\def\qqqquad{\qquad\qquad}

%%%%%%%%%%%%%%%%%%%%%%%%%%%%%%%%%%%%%%%%%%%%%%%%%%%%%%%%%%%%%%%%%%%%%%%%%%%%%%%%
%%%%%%%%%%%%%%%%%%%%%%%%%%%%%%%%%%%%%%%%%%%%%%%%%%%%%%%%%%%%%%%%%%%%%%%%%%%%%%%%

% fill-in-blank question style, found in https://tex.stackexchange.com/a/505089

\usepackage{ifthen}
\usepackage{tocloft}
\usepackage{exercise}
% \usepackage{xcolor}

% Set the Show Answers Boolean
\newboolean{showAns}
\setboolean{showAns}{false}
\newcommand{\showAns}{\setboolean{showAns}{true}}

% The length of the Answer line
\newlength{\answerlength}
\newcommand{\anslen}[1]{\settowidth{\answerlength}{#1}}

% ans command that indicates space for an answer or shows the answer in red
\newcommand{\ans}[1]{\settowidth{\answerlength}{\hspace{2ex}#1\hspace{2ex}}%
    \ifthenelse{\boolean{showAns}}%
        {\textcolor{red}{\underline{\hspace{2ex}#1\hspace{2ex}}}}%
        {\underline{\hspace{\answerlength}}}}%

\newcommand{\details}[1]{\settowidth{\answerlength}{#1}%
    \ifthenelse{\boolean{showAns}}%
        {\begin{tcolorbox}[colback=blue!5!white,colframe=blue!75!black,title=Solution] #1 \end{tcolorbox}}%
        {}}%

% Formatting how multiple choices Questions are formated.
\settasks{label=(\Alph*), label-width=30pt}


% Some commands for the Exercise Question package
\renewcommand{\QuestionNB}{\Large\protect\textcircled{\small\bfseries\arabic{Question}}\ }
\renewcommand{\ExerciseHeader}{} %no header
\renewcommand{\QuestionBefore}{3ex} %Space above each Q
\setlength{\QuestionIndent}{8pt} % Indent after Q number


% To create the list of answers with tocloft...
\newcommand{\listanswername}{Answers}
\newlistof[Question]{answer}{Answers}{\listanswername}

% Creates a TOC for Answers
\newcounter{prevQ}
\newcommand{\answer}[1]{\refstepcounter{answer}%
\ans{#1}%
\ifnum\theQuestion=\theprevQ%
        \addcontentsline{Answers}{answer}{\protect\numberline{}#1}% don't include the Q number
        \else%
        \addcontentsline{Answers}{answer}{\protect\numberline{\theQuestion}#1}%
        \setcounter{prevQ}{\value{Question}}%
        \fi%
        }%

% \hyphenpenalty=10000 % no hyphen
% \exhyphenpenalty=10000 % no hyphen
\sloppy              % hyphen

\newcommand{\HRule}{\rule{\linewidth}{0.5mm}}
\newcommand{\Hrule}{\rule{\linewidth}{0.3mm}}

%tocloft formatting listofanswers
\renewcommand{\cftAnswerstitlefont}{\bfseries\large}
\renewcommand{\cftanswerdotsep}{\cftnodots}
\cftpagenumbersoff{answer}
\addtolength{\cftanswernumwidth}{10pt}

\makeatletter% since there's an at-sign (@) in the command name
\renewcommand{\@maketitle}{%
  \parindent=0pt% don't indent paragraphs in the title block
  \centering
  {\Large \bfseries\textsc{\@title}} \\
  \vspace{5pt}
  {\large \textit{\@author}} \\
  \HRule \\
  \vspace{1em}
}
\makeatother% resets the meaning of the at-sign (@)


% --- Fonts (robust fallback) ---
\IfFontExistsTF{Crimson Pro Light}{
  \setmainfont{Crimson Pro Light}[
    ItalicFont={* Italic},
    BoldFont={Crimson Pro Medium},
    BoldItalicFont={Crimson Pro Medium Italic}]
  \setsansfont{Crimson Pro Light}[
    ItalicFont={* Italic},
    BoldFont={Crimson Pro Medium},
    BoldItalicFont={Crimson Pro Medium Italic}]
}{
  \setmainfont{HaranoAjiMincho}
  \setsansfont{HaranoAjiGothic}
}
\title{Final Exam}
\author{Macroeconomics II \\ Hui-Jun Chen}


% ------------------------------------------------------------
% Explanations for each option (only printed when \showAns is enabled)
% Usage: \ExplainChoices{A}{B}{C}{D}
% ------------------------------------------------------------
\newcommand{\ExplainChoices}[4]{%
\ifthenelse{\boolean{showAns}}{%
\vspace{0.3em}\noindent\textcolor{red}{\textbf{Explanations.}}
\begin{itemize}
\item[(A)] #1
\item[(B)] #2
\item[(C)] #3
\item[(D)] #4
\end{itemize}
}{}%
}
\begin{document}

\maketitle

% Uncomment to display answers + explanations in red
\showAns

\begin{Exercise}
\section*{Problem 1 (Questions 1--20): Inflation, IS--LM, AD--AS}

\Question \textbf{Definition. If $p_t\equiv\log P_t$, then inflation is}
\answer{B}
\begin{tasks}(1)
\task $\pi_t \equiv P_t-P_{t-1}$
\task $\pi_t \equiv p_t-p_{t-1}$
\task $\pi_t \equiv p_{t+1}-p_t$
\task $\pi_t \equiv \log(P_{t-1}/P_t)$
\end{tasks}
\ExplainChoices{Wrong: $P_t-P_{t-1}$ depends on units and is not the standard definition.}{Correct: inflation is the change in the log price level.}{Wrong timing: that is next period’s inflation.}{Wrong sign: this is negative of inflation.}

\Question \textbf{Calculation. Price level goes from 250 to 260 in one year. Approx. inflation is}
\answer{A}
\begin{tasks}(1)
\task $4\%$
\task $10\%$
\task $2.5\%$
\task $-4\%$
\end{tasks}
\ExplainChoices{Correct: $(260-250)/250=0.04$.}{Wrong: would correspond to 25→27.5 or similar.}{Wrong: $(260-250)/260\approx3.85\%$, not 2.5\%.}{Wrong sign.}

\Question \textbf{Fisher relation (ex-ante). The real rate is approximately}
\answer{B}
\begin{tasks}(1)
\task $r_t=i_t+\mathbb{E}_t\pi_{t+1}$
\task $r_t=i_t-\mathbb{E}_t\pi_{t+1}$
\task $r_t=\pi_t-i_t$
\task $r_t=i_t-\pi_t$
\end{tasks}
\ExplainChoices{Wrong sign: expected inflation lowers the real rate.}{Correct: $r\approx i-\mathbb{E}\pi$.}{Wrong sign.}{Wrong: uses current instead of expected inflation.}

\Question \textbf{Calculation. If $i_t=5\%$ and $\mathbb{E}_t\pi_{t+1}=3\%$, then $r_t\approx$}
\answer{B}
\begin{tasks}(1)
\task $8\%$
\task $2\%$
\task $-2\%$
\task $3\%$
\end{tasks}
\ExplainChoices{Wrong: adds instead of subtracts.}{Correct: $5-3=2\%$.}{Wrong sign.}{Wrong: equals expected inflation, not real rate.}

\Question \textbf{Nominal anchor. A nominal anchor is best described as}
\answer{B}
\begin{tasks}(1)
\task a mechanism pinning down real GDP permanently
\task a mechanism pinning down the price level path (or long-run inflation)
\task a mechanism pinning down unemployment at zero
\task a mechanism pinning down money velocity
\end{tasks}
\ExplainChoices{Wrong: real GDP is not pinned by a nominal anchor.}{Correct.}{Wrong: natural unemployment is not zero.}{Wrong: velocity can vary.}

\Question \textbf{Fiscal backing intuition. If expected future primary surpluses fall while nominal debt $B$ is given, then (all else equal)}
\answer{A}
\begin{tasks}(1)
\task $P$ rises so that $B/P$ falls
\task $P$ falls so that $B/P$ rises
\task $P$ is unchanged because $B$ is nominal
\task inflation must be zero
\end{tasks}
\ExplainChoices{Correct: lower backing lowers real value; with nominal $B$, $P$ rises.}{Wrong sign.}{Wrong: $P$ can adjust to revalue nominal claims.}{Wrong.}

\Question \textbf{One-time jump vs persistent inflation. A permanent one-time jump in $P$ implies}
\answer{B}
\begin{tasks}(1)
\task permanently high inflation
\task a one-time burst of inflation, then inflation returns to normal
\task permanently negative inflation
\task inflation is undefined
\end{tasks}
\ExplainChoices{Wrong: persistent inflation requires continuing increases.}{Correct.}{Wrong.}{Wrong.}

\Question \textbf{Money demand. In $M/P=L(Y,i)$, the usual signs are}
\answer{B}
\begin{tasks}(1)
\task $L_Y<0,\\ L_i>0$
\task $L_Y>0,\\ L_i<0$
\task $L_Y>0,\\ L_i>0$
\task $L_Y<0,\\ L_i<0$
\end{tasks}
\ExplainChoices{Wrong: money demand rises with income.}{Correct: higher income raises money demand; higher interest lowers it.}{Wrong: $L_i$ should be negative.}{Wrong.}

\Question \textbf{LM algebra. If $M/P=kY-hi$, then solving for $i$ gives}
\answer{A}
\begin{tasks}(1)
\task $i=\frac{k}{h}Y-\frac{1}{h}\frac{M}{P}$
\task $i=\frac{h}{k}Y-\frac{1}{k}\frac{M}{P}$
\task $i=\frac{1}{h}\frac{M}{P}-\frac{k}{h}Y$
\task $i=kY-h\frac{M}{P}$
\end{tasks}
\ExplainChoices{Correct.}{Wrong: inverted coefficients.}{Wrong sign.}{Wrong.}

\Question \textbf{Numerical LM. Let $M/P=6$, $k=1$, $h=2$, $Y=8$. Then $i$ equals}
\answer{C}
\begin{tasks}(1)
\task $-1$
\task $0$
\task $1$
\task $2$
\end{tasks}
\ExplainChoices{Wrong: compute gives +1.}{Wrong: would require $Y=M/P=6$.}{Correct: $i=0.5\cdot8-0.5\cdot6=1$.}{Wrong.}

\Question \textbf{IS curve sign. In $Y=\bar A-a r$ with $a>0$, higher $r$ implies}
\answer{B}
\begin{tasks}(1)
\task $Y$ rises
\task $Y$ falls
\task $Y$ unchanged
\task $P$ falls mechanically
\end{tasks}
\ExplainChoices{Wrong: higher real rates reduce demand.}{Correct.}{Wrong.}{Wrong: price level is not in IS directly.}

\Question \textbf{Numerical IS. If $Y=100-20r$ and $r$ rises from 0.02 to 0.03, then $\Delta Y=$}
\answer{B}
\begin{tasks}(1)
\task $+0.2$
\task $-0.2$
\task $-2$
\task $-0.4$
\end{tasks}
\ExplainChoices{Wrong sign.}{Correct: $-20\times0.01=-0.2$.}{Wrong magnitude (treating 0.01 as 1).}{Wrong magnitude.}

\Question \textbf{Deriving AD in IS--LM. If $M$ is fixed, a higher price level $P$ typically leads to}
\answer{B}
\begin{tasks}(1)
\task higher $M/P$ and higher $Y$
\task lower $M/P$ and lower $Y$
\task lower $M/P$ and higher $Y$
\task no change in $Y$
\end{tasks}
\ExplainChoices{Wrong: $M/P$ falls when $P$ rises.}{Correct: lower real balances raise $i$ and reduce demand.}{Wrong sign for $Y$.}{Wrong.}

\Question \textbf{Fiscal expansion in IS--LM (short run). A higher $G$ primarily shifts}
\answer{B}
\begin{tasks}(1)
\task LM right
\task IS right
\task LM left
\task IS left
\end{tasks}
\ExplainChoices{Wrong: money market unchanged.}{Correct.}{Wrong.}{Wrong.}

\Question \textbf{Monetary expansion in IS--LM. Higher $M$ (holding $P$ fixed) shifts}
\answer{A}
\begin{tasks}(1)
\task LM right/down
\task IS right
\task LM left/up
\task IS left
\end{tasks}
\ExplainChoices{Correct.}{Wrong: goods-market autonomous demand unchanged.}{Wrong direction.}{Wrong.}

\Question \textbf{AD--AS diagnosis. A negative supply shock typically produces}
\answer{B}
\begin{tasks}(1)
\task $Y\uparrow$ and $P\downarrow$
\task $Y\downarrow$ and $P\uparrow$
\task $Y\uparrow$ and $P\uparrow$
\task $Y\downarrow$ and $P\downarrow$
\end{tasks}
\ExplainChoices{Wrong.}{Correct (stagflation).}{Wrong: demand shock pattern.}{Wrong: that’s contractionary demand shock.}

\Question \textbf{Linear SRAS: if $y=y^{FE}+\alpha(p-p^e)$, then the output gap is}
\answer{A}
\begin{tasks}(1)
\task $x=y-y^{FE}=\alpha(p-p^e)$
\task $x=y+y^{FE}$
\task $x=p-p^e$
\task $x=\alpha(p^e-p)$
\end{tasks}
\ExplainChoices{Correct.}{Wrong.}{Wrong: missing $\alpha$ scaling and units.}{Wrong sign.}

\Question \textbf{Solve AD--AS once. With AD $y=\bar y_D-\phi p$ and SRAS $p=p^e+\lambda(y-y^{FE})$, equilibrium $p$ equals}
\answer{B}
\begin{tasks}(1)
\task $p=p^e+\lambda(\bar y_D-y^{FE})$
\task $p=\dfrac{p^e+\lambda(\bar y_D-y^{FE})}{1+\lambda\phi}$
\task $p=\dfrac{p^e}{1+\lambda\phi}$
\task $p=\dfrac{\bar y_D-y^{FE}}{1+\lambda\phi}$
\end{tasks}
\ExplainChoices{Wrong: forgot feedback through AD slope $\phi$.}{Correct: substitute AD into SRAS and solve for $p$.}{Wrong: ignores demand term.}{Wrong: ignores $p^e$.}

\Question \textbf{Okun numerical. If $u-u^n=-0.5 x$ and $x=2\%$, then $u-u^n$ is}
\answer{B}
\begin{tasks}(1)
\task $+1\%$
\task $-1\%$
\task $+4\%$
\task $-4\%$
\end{tasks}
\ExplainChoices{Wrong sign.}{Correct: $-0.5\times2\%=-1\%$.}{Wrong magnitude.}{Wrong magnitude.}

\Question \textbf{IS--MP--PC model blocks are}
\answer{B}
\begin{tasks}(1)
\task IS, LM, and Phillips curve
\task IS, monetary policy rule, and Phillips curve
\task money supply rule, SRAS, and LRAS
\task government budget constraint, money demand, and NKPC
\end{tasks}
\ExplainChoices{Wrong: LM is replaced by a policy rule in IS–MP–PC.}{Correct.}{Wrong.}{Wrong.}

\vspace{0.6em}\Hrule\vspace{0.8em}
\section*{Problem 2 (Questions 21--40): IS--MP--PC, Taylor principle, NK}

\Question \textbf{Whelan-style Phillips curve: $\pi_t=\pi_t^e+\gamma(y_t-y_t^*)+\varepsilon_t^\pi$. If $\pi_t^e$ rises holding gap fixed, then}
\answer{A}
\begin{tasks}(1)
\task $\pi_t$ rises one-for-one
\task $\pi_t$ falls one-for-one
\task $\pi_t$ unchanged
\task the slope $\gamma$ changes
\end{tasks}
\ExplainChoices{Correct: intercept shifts up.}{Wrong.}{Wrong.}{Wrong: slope unchanged.}

\Question \textbf{Baseline MP rule: $i_t=r^*+\pi^*+\beta_\pi(\pi_t-\pi^*)$. The Taylor principle typically requires}
\answer{C}
\begin{tasks}(1)
\task $\beta_\pi<1$
\task $\beta_\pi=1$
\task $\beta_\pi>1$
\task $\beta_\pi<0$
\end{tasks}
\ExplainChoices{Wrong.}{Boundary (real rate doesn’t respond to inflation deviations).}{Correct.}{Wrong.}

\Question \textbf{IS--MP slope sign. In IS--MP, output falls with inflation (downward IS--MP in $(y,\pi)$) when}
\answer{C}
\begin{tasks}(1)
\task $\beta_\pi<1$
\task $\beta_\pi=0$
\task $\beta_\pi>1$
\task $\gamma=0$
\end{tasks}
\ExplainChoices{Wrong: then IS--MP slopes up (destabilizing).}{Not the key condition.}{Correct: real rate rises when inflation rises.}{Not relevant for IS--MP slope.}

\Question \textbf{Define $\theta=\dfrac{1}{1+\alpha\gamma(\beta_\pi-1)}$ with $\alpha,\gamma>0$. If $\beta_\pi>1$, then}
\answer{B}
\begin{tasks}(1)
\task $\theta>1$
\task $0<\theta<1$
\task $\theta=1$
\task $\theta<0$
\end{tasks}
\ExplainChoices{Wrong: denominator exceeds 1 so cannot exceed 1.}{Correct.}{Wrong: occurs when $\beta_\pi=1$ (or $\alpha\gamma=0$).}{Wrong.}

\Question \textbf{Numerical $\theta$. Let $\alpha=1$, $\gamma=0.5$, $\beta_\pi=3$. Then $\theta=$}
\answer{B}
\begin{tasks}(1)
\task $0.4$
\task $0.5$
\task $2$
\task $-1$
\end{tasks}
\ExplainChoices{Wrong: used $\beta_\pi$ instead of $(\beta_\pi-1)$.}{Correct: $1/(1+0.5\cdot2)=1/2=0.5$.}{Wrong.}{Wrong.}

\Question \textbf{Interpretation. A smaller $\theta$ means inflation is}
\answer{B}
\begin{tasks}(1)
\task more sensitive to expectations
\task more tightly anchored to the target
\task more sensitive to shocks
\task independent of policy
\end{tasks}
\ExplainChoices{Wrong: smaller $\theta$ reduces weight on expectations/shocks.}{Correct.}{Wrong.}{Wrong.}

\Question \textbf{Demand shock in IS--MP--PC (one-period). A positive demand shock shifts IS--MP right and typically leads to}
\answer{A}
\begin{tasks}(1)
\task $y\uparrow$ and $\pi\uparrow$
\task $y\downarrow$ and $\pi\uparrow$
\task $y\uparrow$ and $\pi\downarrow$
\task $y\downarrow$ and $\pi\downarrow$
\end{tasks}
\ExplainChoices{Correct.}{Cost-push pattern.}{Wrong.}{Wrong.}

\Question \textbf{Cost-push shock ($\varepsilon^\pi>0$) typically leads to}
\answer{B}
\begin{tasks}(1)
\task $y\uparrow$ and $\pi\uparrow$
\task $y\downarrow$ and $\pi\uparrow$
\task $y\uparrow$ and $\pi\downarrow$
\task $y\downarrow$ and $\pi\downarrow$
\end{tasks}
\ExplainChoices{Wrong.}{Correct (stagflation pressure).}{Wrong.}{Wrong.}

\Question \textbf{Adaptive expectations assumption is}
\answer{B}
\begin{tasks}(1)
\task $\pi_t^e=\pi^*$
\task $\pi_t^e=\pi_{t-1}$
\task $\pi_t^e=\mathbb{E}_t\pi_{t+1}$
\task $\pi_t^e=0$ always
\end{tasks}
\ExplainChoices{Wrong (anchored).}{Correct.}{Forward-looking.}{Wrong.}

\Question \textbf{Two-period NK IS equation is}
\answer{A}
\begin{tasks}(1)
\task $x_0=\mathbb{E}_0x_1-\frac{1}{\sigma}(i_0-\mathbb{E}_0\pi_1-r_0^n)$
\task $x_0=\mathbb{E}_0x_1+\frac{1}{\sigma}(i_0-\mathbb{E}_0\pi_1-r_0^n)$
\task $x_0=\frac{1}{\sigma}(i_0-\mathbb{E}_0\pi_1-r_0^n)$
\task $x_0=-\sigma(i_0-\mathbb{E}_0\pi_1-r_0^n)$
\end{tasks}
\ExplainChoices{Correct.}{Wrong sign.}{Missing expectation term and wrong sign.}{Wrong scaling.}

\Question \textbf{Two-period NKPC is}
\answer{B}
\begin{tasks}(1)
\task $\pi_0=\pi_{-1}+\kappa x_0+u_0$
\task $\pi_0=\beta\mathbb{E}_0\pi_1+\kappa x_0+u_0$
\task $\pi_0=\kappa x_0+u_0$ always
\task $\pi_0=\beta\pi_0+\kappa x_0$
\end{tasks}
\ExplainChoices{Backward-looking.}{Correct.}{Only if $\mathbb{E}_0\pi_1=0$.}{Wrong.}

\Question \textbf{Calculation (NKPC). Let $\beta=0.99$, $\mathbb{E}_0\pi_1=0$, $\kappa=0.5$, $u_0=0$, and $x_0=2\%$. Then $\pi_0=$}
\answer{A}
\begin{tasks}(1)
\task $1\%$
\task $0.5\%$
\task $2\%$
\task $0\%$
\end{tasks}
\ExplainChoices{Correct: $\pi_0=0.5\times2\%=1\%$.}{Wrong: would be $x_0=1\%$.}{Wrong: would require $\kappa=1$.}{Wrong.}

\Question \textbf{Calculation (Taylor rule). If $i_0=\bar i+\phi_\pi\pi_0$ with $\bar i=1\%$, $\phi_\pi=1.5$, and $\pi_0=2\%$, then $i_0=$}
\answer{C}
\begin{tasks}(1)
\task $2\%$
\task $3\%$
\task $4\%$
\task $5\%$
\end{tasks}
\ExplainChoices{Wrong: missing reaction term.}{Wrong: used $\phi_\pi=1$ effectively.}{Correct: $1+1.5\times2=4\%$.}{Wrong.}

\Question \textbf{Calculation (Fisher). If $i_0=4\%$ and $\mathbb{E}_0\pi_1=1\%$, then $r_0\approx$}
\answer{A}
\begin{tasks}(1)
\task $3\%$
\task $5\%$
\task $-3\%$
\task $1\%$
\end{tasks}
\ExplainChoices{Correct.}{Wrong.}{Wrong sign.}{Wrong.}

\Question \textbf{ZLB. The zero lower bound is the constraint}
\answer{B}
\begin{tasks}(1)
\task $r_0\ge 0$
\task $i_0\ge 0$
\task $\pi_0\ge 0$
\task $x_0\ge 0$
\end{tasks}
\ExplainChoices{Wrong: real rates can be negative.}{Correct.}{Wrong.}{Wrong.}

\Question \textbf{ZLB implication. If $i_0=0$ and expected inflation falls, then the real rate}
\answer{B}
\begin{tasks}(1)
\task falls
\task rises
\task is unchanged
\task becomes negative
\end{tasks}
\ExplainChoices{Wrong: $r=i-\mathbb{E}\pi$.}{Correct: lower expected inflation raises $r$.}{Wrong.}{Wrong: sign depends on inflation expectation.}

\Question \textbf{Calculation (IS at ZLB). Suppose $x_1=0$, $\sigma=1$, $r_0^n=0\%$, $i_0=0$, and $\mathbb{E}_0\pi_1=-1\%$. Then $x_0=$}
\answer{C}
\begin{tasks}(1)
\task $0$
\task $+1\%$
\task $-1\%$
\task $+2\%$
\end{tasks}
\ExplainChoices{Wrong: real rate gap is positive, so gap is negative.}{Wrong sign.}{Correct: $i-\mathbb{E}\pi-r^n=0-(-1)-0=1\%$, so $x_0=-(1\%)=-1\%$.}{Wrong magnitude.}

\Question \textbf{Calculation (IS–MP). If IS is $y=y^*-\alpha(i-\pi-r^*)$ and MP is $i=r^*+\pi^*+\beta_\pi(\pi-\pi^*)$, then the coefficient on $(\pi-\pi^*)$ in IS--MP is}
\answer{A}
\begin{tasks}(1)
\task $-\alpha(\beta_\pi-1)$
\task $+\alpha(\beta_\pi-1)$
\task $-\alpha\beta_\pi$
\task $0$
\end{tasks}
\ExplainChoices{Correct: substitute MP into IS and simplify $i-\pi-r^*=(\beta_\pi-1)(\pi-\pi^*)$.}{Wrong sign.}{Wrong: misses the $-\pi$ term when forming the real rate gap.}{Wrong: only if $\beta_\pi=1$.}

\Question \textbf{Calculation (policy aggressiveness). If $\alpha=1$ and $\gamma=0.5$, what value of $\beta_\pi$ makes the boundary $1-\frac{1}{\alpha\gamma}$ equal to $-1$?}
\answer{C}
\begin{tasks}(1)
\task $\beta_\pi=0$
\task $\beta_\pi=1$
\task $\beta_\pi$ does not enter this boundary; the boundary is $-1$ given $\alpha\gamma=0.5$.
\task $\beta_\pi=2$
\end{tasks}
\ExplainChoices{Wrong: boundary depends only on $\alpha\gamma$.}{Wrong.}{Correct: $1-1/(0.5)= -1$ regardless of $\beta_\pi$.}{Wrong.}

\Question \textbf{Concept. In the NK model, forward-looking behavior is primarily captured by}
\answer{B}
\begin{tasks}(1)
\task $\pi_t=\pi_{t-1}+\kappa x_t$
\task $x_t=\mathbb{E}_t x_{t+1}-\frac{1}{\sigma}(i_t-\mathbb{E}_t\pi_{t+1}-r_t^n)$ and $\pi_t=\beta\mathbb{E}_t\pi_{t+1}+\kappa x_t+u_t$
\task $M/P=kY-hi$
\task $y=y^{FE}+\alpha(p-p^e)$
\end{tasks}
\ExplainChoices{Wrong: backward-looking.}{Correct: expectations of next-period variables enter both IS and NKPC.}{Wrong: money demand is IS–LM.}{Wrong: SRAS is a reduced-form level relation.}

\end{Exercise}
\end{document}
